\begin{tabular}{@{} >{\small}l@{} >{\raggedright\arraybackslash\small}p{0.25\textwidth} >{\raggedright\arraybackslash\small}p{0.57\textwidth}@{}}
    \multicolumn{3}{@{}l}{\textbf { Process Quality} \small {\textit{(relates to the democratic infrastructure)}}}  \\

    \midrule
    
    & \textit{The extent to which ...} & \\

    Representation
        & key decisions are representative of the constituent population.
        & To what extent: \b{1} is there sufficient representation at critical parts of the process, including (a) proposing decisions, and (b) making ultimate decisions? \b{2} are there barriers leading to bias in representation? \\

    Informedness
        & critical information is taken into account for decision-making.
        & To what extent: \b{1} is critical context incorporated into decision-making about tradeoffs and consequences of different decisions? \b{2} is this sourced from (a) domain expertise, (b) the existing authorities, who may have extensive context, (c) a broad diversity of constituents, (d) the most impacted stakeholders, and (e) the powerful stakeholders, whose incentives are critical to having the decision ``stick''? \\

    Deliberation
        & decisions are considered and deliberative (rather than superficial and reactive).
        & To what extent are those involved: \b{1} able to (and supported to) move from shallower to deeper goals and values? \b{2} able to (and supported to) collaborate where necessary? \b{3} able to address issues within the available time?
        \\

    Substantiveness\ \ \ \ 
        & decisions are substantive (e.g., actionable, consequential) rather than nonsubstantive (e.g., vague, simplistic, inconsequential). 
        & To what extent: \b{1} is the decision directly actionable and implementable? \b{2} does the decision meaningfully address the issues? \b{3} does the decision grapple with the necessary levels of complexity? \b{4} is uncertainty appropriately managed and accounted for? \b{5} are risks to implementability accounted for? \\

    Robustness
        & the process is robust to suboptimal conditions or adversarial or strategic behavior.
        & To what extent is the process or system vulnerable to: \b{1} suboptimal conditions or broken assumptions? (e.g., low turnout, larger power asymmetries) \b{2} strategic behavior and manipulation? \b{3} false claims? (e.g., of manipulation) \\

    Legibility
        & the processes and decisions are accessible, understandable, and verifiable %
        & To what extent is information (a) accessible, (b) understandable, (c) verifiable about the: \b{1} processes/systems used to make decisions? \b{2} the execution of these processes? \b{3}~decisions being made \b{4} reasons and inputs feeding into decisions? \\[10pt]

    \multicolumn{3}{@{}l}{\textbf { Delegation} \small{\textit{(relates to the unilateral or existing authorities)}}}   \\

    \midrule
    
    Integration
        & the authority integrates the democratic process into its operations.
        & To what extent is the authority structuring its internal communications and operations to effectively: \b{1} provide critical context to democratic process / system? \b{2} integrate democratic process outputs in its actions? \b{3} trigger democratic processes when/if required?
        \\

    Ability to bind
        & the authority is able to technically and legally bind itself to democratic decisions.
        & To what extent can the unilateral authority bind itself to acting in accordance with the democratic decision: \b{1} technically? \b{2} legally? (e.g., has developed the needed technical and/or legal infrastructure for binding)
        \\
        
    Commitment
        & the unilateral authority commits to acting in accordance with the democratic decision.
        & To what extent has the unilateral authority committed to acting in accordance with the democratic decision: \b{1} internally? \b{2} privately? \b{3} publicly?  (regardless of their ability to bind) \\[10pt]

    \multicolumn{3}{@{}l}{\textbf { Trust}  \small{\textit{(relates to the external conditions)}}}\\
    \midrule
    Awareness
        & the relevant public is aware of the democratic process.
        & To what extent is the relevant public aware: \b{1} that the democratic system exists? \b{2} how it works? \b{3} what it is being used for? \b{4} how they can be involved? \\

    Participation
        & the relevant public is willing to participate in the process.
        & To what extent is the relevant public: \b{1} willing to participate? \b{2} able to participate? \b{3} appropriately compensated for participating? \b{4} actually participating? \\
        
    Accountability
        & there are external watchdogs and accountability structures monitoring the execution of the democratic process and the implementation of its outputs.
        & To what extent are: \b{1} there well understood lines of oversight and accountability? \b{2} sufficiently influential/powerful organizations focused on holding authorities to their promised levels of democratic involvement? \b{3} authorities and democratic systems responsive to such accountability mechanisms? \\
        
    Buy-in
        & the relevant public and key stakeholders buy-in to the process and its legitimacy.
        & To what extent are the relevant public and key stakeholders accepting of the legitimacy of: \b{1} the system/process? \b{2} of the decision?\\
        
    
    \cmidrule[\heavyrulewidth]{1-3} %
\end{tabular}
